\documentclass[11pt,a4paper]{article}
\usepackage[utf8]{inputenc}
\usepackage[T1]{fontenc}
\usepackage{geometry}
\geometry{margin=2.5cm}
\usepackage{hyperref}
\usepackage{listings}
\usepackage{xcolor}
\usepackage{enumitem}
\usepackage{booktabs}
\usepackage{graphicx}
\usepackage{tikz}
\usetikzlibrary{arrows.meta,positioning,shapes.geometric}

\lstset{
  basicstyle=\ttfamily\small,
  breaklines=true,
  frame=single,
  backgroundcolor=\color{gray!10},
  keywordstyle=\color{blue!70!black},
  commentstyle=\color{green!50!black},
  stringstyle=\color{red!60!black},
  showstringspaces=false,
  tabsize=2
}

\title{Time as a First-Class Control Surface in Shadow:\\Methodology, Capabilities, and Implementation Challenges}
\author{Shadow Enhancement Project}
\date{April 2026}

\begin{document}
\maketitle

\begin{abstract}
Traditional uses of Shadow treat simulated time primarily as an internal
scheduler variable. Our refactoring elevates time to a \emph{first-class
control surface}: external tools and interactive users can now pause the
system at globally consistent boundaries, step the execution in a controlled
fashion, drive the simulator to a target logical time, and trigger
checkpoint/restore transitions. The resulting interface is motivated not by
operator convenience alone, but by higher-level systems goals such as
time-travel debugging, snapshot-based fuzzing, deterministic replay, and
stateful experiment orchestration.

The key methodological choice is to align every control action with
Shadow's runahead window boundaries. This yields a well-defined global cut
over a parallel discrete-event execution and avoids unsafe interventions in
the middle of host execution, shim IPC, or packet transfer. Building such a
time-control substrate is non-trivial because Shadow combines
(i) parallel event-driven scheduling,
(ii) real OS processes executing unmodified binaries,
(iii) Rust-side simulator state that is rich but not uniformly serializable,
and (iv) external stateful dependencies such as databases or ledgers that
live outside Shadow's address space.

This document summarizes the core methodology, the capabilities now exposed
by the implementation, the decomposition of the design into control-plane and
state-plane components, and the principal challenges encountered during
implementation, together with the mechanisms used to overcome them.
\end{abstract}

\tableofcontents
\newpage

\section{Introduction}

The central thesis of this work is that \emph{time control} should be a
native systems abstraction in Shadow, rather than an ad hoc debugging aid.
For many workloads executed in Shadow---especially consensus clients,
distributed databases, network services, and event-driven state machines---the
most interesting failures are localized not merely by host identity or packet
ordering, but by \emph{logical time}. Reaching, freezing, replaying, and
restoring a particular temporal region is therefore a first-class requirement
for modern systems experimentation.

This observation motivates a shift from the legacy view of Shadow's control
logic as a local convenience embedded inside \texttt{manager.rs} toward a
more principled design in which simulated time is exposed as an explicit and
programmable interface. Concretely, the current implementation provides a
time-control substrate with four primary classes of capability:

\begin{enumerate}[leftmargin=1.7em]
  \item \textbf{Boundary-aligned pause/resume}: freeze execution at a
  globally consistent logical-time boundary and resume without perturbing the
  scheduler.
  \item \textbf{Temporal stepping and run-until}: advance one runahead window
  at a time, or execute until a specified simulated-time target.
  \item \textbf{Checkpoint}: materialize a restorable temporal cut, including
  simulator metadata, event representations, shared-memory artifacts, and
  process-level checkpoint hooks.
  \item \textbf{Restore}: re-enter the simulation from a prior temporal cut,
  enabling replay-style experimentation and the basis for time-travel
  debugging.
\end{enumerate}

These capabilities are not ends in themselves. They are intended to support
systematic analyses that are otherwise expensive or fragile in real-time
testbeds, including (i) deterministic reproduction of narrow fault windows,
(ii) debugger attachment to the precise processes scheduled in the next
logical-time slice, (iii) branch-and-replay exploration for fuzzing and fault
injection, and (iv) coordinated control over Shadow and external persistent
state.

\section{Methodological Core}

\subsection{Time Control as a Boundary-Synchronous Abstraction}

Our methodology is based on a simple but fundamental principle:
\emph{time-control decisions should only take effect at a globally safe
simulation boundary}. In Shadow, the natural boundary is the runahead window
boundary already induced by the parallel scheduler. At such a boundary,
the simulator has completed all causally ready work within the current
window and is about to determine the next execution frontier. This point is
therefore uniquely suitable for intervention.

This choice has three methodological consequences.

\paragraph{Consistency.}
Pausing at a window boundary yields a consistent global cut. It avoids stopping
in the middle of a host's internal execution, a partially serviced syscall,
or an in-flight Shadow--shim exchange.

\paragraph{Determinism preservation.}
Stepping and run-until semantics are expressed in terms of window transitions
rather than wall-clock preemption. This preserves the existing event ordering
discipline of the simulator and avoids introducing extra nondeterministic
interleavings.

\paragraph{Composability.}
Checkpoint and restore become naturally aligned with the same boundary notion:
the checkpoint is taken at a safe cut, and restore is defined as restarting
the simulation from a previously recorded cut and re-establishing the state
required to continue from that point.

\subsection{Why This Matters}

The boundary-synchronous model is particularly important when the goal is not
merely to ``pause the simulator'', but to \emph{treat time as an experimental
dimension}. This is the enabling condition for several higher-level workflows:

\begin{itemize}[leftmargin=1.7em]
  \item \textbf{Time-travel debugging.} A developer can advance to a target
  interval, checkpoint it, inspect the processes that will execute next, and
  replay or restore the same region under different debugging conditions.
  \item \textbf{Snapshot-based fuzzing.} A fuzzer can checkpoint a promising
  precondition, fork exploration from that cut, and repeatedly restore to the
  same logical-time state while varying inputs or perturbations.
  \item \textbf{Reproducible fault localization.} Failures observed at a given
  logical time can be revisited using deterministic run-until or restore,
  instead of re-running the full experiment from scratch without guarantees of
  precise temporal alignment.
  \item \textbf{Stateful orchestration.} External controllers can coordinate
  Shadow with databases, ledgers, or file-backed state that cannot be safely
  reset by the simulator alone.
\end{itemize}

\section{System Design}

\subsection{Control Plane}

The control plane is factored around a trait-based abstraction:

\begin{lstlisting}[language=Rust]
pub trait TimeController: Send + Sync {
    fn on_simulation_start(&self);
    fn on_window_boundary(
        &self,
        ctx: &WindowBoundaryContext,
        print_info: PrintNextWindowInfoFn<'_>,
    ) -> ControlDecision;
    fn on_simulation_end(&self);
}
\end{lstlisting}

This interface separates \emph{control policy} from the simulation kernel.
The manager no longer embeds a monolithic run-control mechanism. Instead, it
invokes a controller at three lifecycle hooks: simulation start, every window
boundary, and simulation end. The important consequence is architectural:
time control is now a pluggable subsystem rather than a special-case block of
logic entangled with the scheduler.

\subsection{Decision Semantics}

The controller returns a \texttt{ControlDecision} interpreted by the manager.
The current design supports the following decision space:

\begin{itemize}[leftmargin=1.7em]
  \item \texttt{Continue}: keep executing the next window.
  \item \texttt{PauseAtBoundary}: remain paused until the front-end explicitly
  resumes.
  \item \texttt{Restart}: restart from logical time \(t=0\), optionally with a
  target run-until time.
  \item \texttt{CheckpointNow}: materialize a checkpoint at the current cut.
  \item \texttt{RestoreCheckpoint}: leave the current run and re-enter from a
  named checkpoint.
\end{itemize}

Equally important is the result returned by \texttt{Manager::run()}:

\begin{lstlisting}[language=Rust]
pub enum SimulationRunResult {
    Completed { num_plugin_errors: u32 },
    RestartRequested {
        run_until_ns: Option<u64>,
        source: RestartSource,
    },
    RestoreRequested { label: String },
}
\end{lstlisting}

This replaces the earlier error-based restart path and makes temporal control
an explicit part of the simulator's normal control flow. The outer loop in
\texttt{shadow.rs} now interprets restart and restore requests as ordinary
state transitions rather than exceptional failures.

\subsection{Front-End Controllers}

The current system includes three concrete controllers.

\paragraph{\texttt{NoopController}.}
This controller always returns \texttt{Continue}. It preserves the legacy
behavior when time control is disabled.

\paragraph{\texttt{InteractiveController}.}
This controller exposes a debugger-like stdin interface. It supports pausing
at the next boundary, continuing indefinitely, continuing for \(N\) simulated
seconds, stepping exactly one runahead window, requesting restart, printing
next-window diagnostics, and emitting attach information for native PIDs.
The semantics are intentionally \emph{soft}: when paused, the simulator blocks
on a condition variable at a boundary rather than attempting in-window
preemption.

\paragraph{\texttt{SocketController}.}
This controller exposes time control to external orchestrators over a Unix
domain socket using newline-delimited JSON. It is designed for workflows in
which Shadow alone is not the whole system under experiment. The socket
front-end supports external pause/continue, bounded execution via
\texttt{continue\_for}, restart/replay, checkpoint, restore, and lightweight
status queries. Critically, the socket controller uses a condition-variable
handshake to ensure that commands are only consumed when the simulator is
parked at a boundary.

\section{Operational Capabilities}

Table~\ref{tab:capabilities} summarizes the practical temporal capabilities
now exposed by the implementation.

\begin{table}[h]
\centering
\begin{tabular}{@{}p{3.1cm}p{3.2cm}p{7.1cm}@{}}
\toprule
\textbf{Capability} & \textbf{Front-end} & \textbf{Purpose and meaning} \\
\midrule
Pause at next boundary & stdin, socket & Freeze the execution at a globally safe logical-time cut for inspection, intervention, or orchestration. \\
Resume / continue & stdin, socket & Continue normal execution without changing the causal schedule. \\
Run until target time & stdin (\texttt{cN}), socket (\texttt{continue\_for}) & Reach a target logical-time region quickly while preserving deterministic window-boundary semantics. \\
Single-window step & stdin (\texttt{n}) & Advance exactly one runahead window, which is the minimal scheduler-aligned stepping unit. \\
Restart / replay & stdin, socket & Re-run from \(t=0\); external mode is intended for coordinated reset of persistent state outside Shadow. \\
Checkpoint & socket / control plane & Materialize a boundary-aligned temporal cut that can later be restored or used as a branch point for exploration. \\
Restore & socket / control plane & Re-enter the experiment from a previously saved cut, enabling temporal rewinding and branching. \\
Next-window inspection & stdin & Reveal which hosts and native processes are scheduled next, providing a bridge from logical time to concrete OS processes. \\
\bottomrule
\end{tabular}
\caption{Time-control capabilities exposed by the current design.}
\label{tab:capabilities}
\end{table}

The most important point is not the command vocabulary itself, but the
semantic contract behind it: every operation is aligned with a scheduler-safe
boundary. This makes the interface suitable for rigorous experimentation
instead of merely ad hoc pausing.

\section{Stateful Time Control: Checkpoint and Restore}

\subsection{Design Rationale}

Pause and step are useful, but they are not sufficient for temporal
branching. Time-travel debugging and snapshot-based fuzzing require the system
to revisit a prior execution state without replaying the entire prefix from
scratch. In Shadow, however, such a capability is substantially harder than
in a pure simulator because the complete execution state is split across
multiple layers:

\begin{enumerate}[leftmargin=1.7em]
  \item \textbf{Rust-side simulator state}: event queues, runahead metadata,
  host-local counters, random state, and process/thread metadata.
  \item \textbf{Managed native processes}: unmodified binaries launched by
  Shadow and executing as real OS processes.
  \item \textbf{Shadow--shim IPC state}: shared-memory regions that mediate
  communication between Shadow and the managed processes.
  \item \textbf{External persistent state}: databases, ledgers, and files that
  are not owned by Shadow.
\end{enumerate}

Accordingly, our checkpoint/restore design is decomposed into cooperating
mechanisms rather than a single monolithic snapshot.

\subsection{Snapshot Decomposition}

The checkpoint module introduces serializable snapshot types that mirror the
live structures while remaining \texttt{serde}-compatible. The top-level
\texttt{SimulationCheckpoint} stores the simulated timestamp, runahead
metadata, host-level checkpoint records, shared-memory backup directories, and
CRIU-related paths. Host-level snapshots record event-queue contents, local
counters, PRNG state, and process metadata. Event-level snapshots distinguish
between packet events and local task events.

This decomposition is methodologically important: the restore path does not
attempt to serialize opaque Rust objects directly. Instead, it defines an
explicit state representation whose purpose is to be persisted, validated, and
reconstructed.

\subsection{Task Reification}

One of the hardest parts of simulator-state checkpointing is that many local
events in Shadow are represented by \texttt{TaskRef}, whose payload is an
\texttt{Arc<dyn Fn(\&Host)>}. Such closures are fundamentally not serializable.
The implementation therefore introduces \texttt{TaskDescriptor}, a parallel
representation that records the \emph{semantic identity} of reconstructable
tasks, such as process resume, shutdown, timer expiration, relay forwarding,
or exec continuation. During event-queue serialization, local tasks are
converted into descriptors; during restore, reconstructable descriptors are
mapped back into live \texttt{TaskRef} closures.

This mechanism turns task serialization from an impossible generic problem
into a controlled, explicitly enumerated one. It also makes limitations
observable: tasks that are opaque or semantically incomplete are marked as
such instead of being silently mishandled.

\subsection{Packet Reification}

Packet state is similarly reified into serializable snapshots containing
protocol headers, payload bytes, and priority metadata. The current
implementation supports TCP and UDP packet reconstruction through public packet
constructors. Some internal TCP representation details, such as structures not
publicly reconstructable from library interfaces, are handled conservatively.
This is a recurring theme in the design: where Shadow's live state depends on
non-public or non-serializable internal forms, the implementation introduces
explicit snapshot structures and rebuild logic instead of relying on raw
memory serialization.

\subsection{Shared Memory and CRIU}

Managed processes and the simulator communicate through shared-memory-backed
IPC channels. A restorable cut therefore requires more than event-queue
metadata: it must preserve the shared-memory artifacts that both sides expect.
The implementation addresses this by backing up the relevant \texttt{/dev/shm}
files into the checkpoint directory and restoring them before the simulator
re-enters the run loop. In parallel, CRIU integration is used to checkpoint
and restore managed OS processes when process images are available.

The key systems insight here is that ``restore'' is not a single action. It is
an orchestrated sequence:

\begin{enumerate}[leftmargin=1.7em]
  \item load checkpoint metadata;
  \item restore shared-memory artifacts;
  \item restore managed processes through CRIU where applicable;
  \item re-enter Shadow's top-level control loop;
  \item re-create simulator structures for continued execution.
\end{enumerate}

\subsection{External Orchestration}

Checkpoint/restore in Shadow cannot be understood correctly without external
state. A simulation may interact with persistent databases or file-backed
application state that lies entirely outside the simulator's ownership.
Accordingly, the current design distinguishes \emph{internal restart} from
\emph{external restart/replay}. The socket-based control path is the preferred
mechanism for stateful workloads because it lets an external orchestrator reset
the database or restore the external filesystem before asking Shadow to
restart or restore.

This division of responsibility is essential for correctness. A simulator can
only guarantee restoration of the state it controls. External orchestration is
therefore not an implementation convenience, but part of the semantic
contract of stateful time control.

\section{Implementation Summary}

\subsection{Integration into the Main Loop}

The refactoring removes the previous monolithic run-control code from
\texttt{manager.rs} and replaces it with a single controller callback at each
window boundary. The manager constructs a \texttt{WindowBoundaryContext}
containing the current simulated time, next-event information, and the window
range. The controller consumes this context and returns a
\texttt{ControlDecision}. The manager then either proceeds, pauses, performs a
checkpoint, or exits the current run with a structured restart/restore
request.

At the outermost layer, \texttt{shadow.rs} hosts a loop over simulation runs.
This loop is now responsible for interpreting \texttt{SimulationRunResult},
invoking restore orchestration when required, and re-entering the simulation
cleanly. This design is significantly clearer than the older error-based
mechanism and better matches the fact that restart and restore are
first-class control transitions.

\subsection{Socket-Level Synchronization}

The socket controller implements a subtle but important synchronization
protocol. Commands received from a client are not applied immediately. Instead,
the controller waits until the simulation thread has parked at a window
boundary, marks the pending control decision, and only then allows the
simulation to consume it. Bounded execution (\texttt{continue\_for}) is
implemented by converting a relative simulated-time duration into an absolute
deadline tracked by the controller. The response to the client is delayed until
the simulator has again reached a boundary after satisfying the request.

This design solves a classic control-plane race: without such a handshake, an
external tool could issue a command while the simulator is between boundaries,
leading either to imprecise response timing or to unsafe assumptions about the
state at which the command took effect.

\subsection{POC Workflow}

The current codebase includes a proof-of-concept orchestrator and a simple
file-backed counter application. The workflow is representative:

\begin{enumerate}[leftmargin=1.7em]
  \item launch Shadow with \texttt{SHADOW\_CONTROL\_SOCKET};
  \item execute until a target logical time;
  \item issue a checkpoint request and back up the external database file;
  \item continue execution to a later time;
  \item restore the external database and issue a restore request;
  \item continue from the restored cut and verify that the application-visible
  state has been rewound.
\end{enumerate}

Conceptually, this POC demonstrates the intended composition of temporal
control over both Shadow-managed state and external persistent state.

\section{Why the Implementation is Non-Trivial}

The implementation required solving several non-trivial systems problems. We
highlight the most important ones here.

\subsection{Challenge 1: Defining a Safe Temporal Intervention Point}

\textbf{Problem.}
In a parallel discrete-event simulator, naively pausing ``right now'' is not a
well-defined operation. Hosts may be in the middle of processing local work,
syscall emulation may be in progress, and cross-host effects may not yet be
globally visible.

\textbf{Resolution.}
We align all control actions with runahead window boundaries. This yields a
global cut that is both scheduler-safe and semantically meaningful. The same
choice also gives stepping a natural unit: one window.

\subsection{Challenge 2: Serializing Local Events that are Implemented as Closures}

\textbf{Problem.}
Shadow's local event system relies on task closures, which cannot be
serialized directly. This blocks straightforward checkpointing of event queues.

\textbf{Resolution.}
We introduced \texttt{TaskDescriptor} as a semantic replacement for closures.
Checkpointing records descriptors instead of trying to serialize executable
code. Restore reconstructs live tasks only for descriptors with an explicit
rebuild path, while conservatively skipping opaque tasks with warnings.

\subsection{Challenge 3: Reconciling Simulator State with Real Process State}

\textbf{Problem.}
Shadow is not a pure simulator; it runs real OS processes. Checkpointing only
the Rust-side scheduler state is insufficient because the application state
also resides in native processes.

\textbf{Resolution.}
We decomposed checkpoint/restore into simulator metadata snapshots plus CRIU
integration for managed processes. This separation respects the fact that
process state and simulator state live in different abstraction layers and
must be restored through different mechanisms.

\subsection{Challenge 4: Preserving Shadow--Shim IPC Across Restore}

\textbf{Problem.}
Even if process state is restorable, the communication channels between Shadow
and the shim depend on shared-memory objects and process identity. A restart of
Shadow without rebuilding this substrate would leave the restored processes
disconnected from the new simulator instance.

\textbf{Resolution.}
The implementation explicitly backs up shared-memory artifacts, restores them
before re-entry, and threads the restore path through top-level orchestration
so that the simulator can rebuild the environment expected by managed
processes. This is also why restore is handled outside the core scheduler loop
rather than as a local in-loop mutation.

\subsection{Challenge 5: External Stateful Dependencies}

\textbf{Problem.}
Realistic workloads often depend on state that Shadow does not own: databases,
blockchain data directories, or other persistent services. Restarting or
restoring only the simulator would not produce a semantically valid replay.

\textbf{Resolution.}
We made the distinction between \emph{internal} and \emph{external} restart
explicit in the API through \texttt{RestartSource}. Internal restart emits a
warning that only Shadow's internal state is reset. External restart and
restore are exposed through the socket controller so that an orchestrator can
reset or restore persistent state in lockstep with Shadow.

\subsection{Challenge 6: Giving Precise Semantics to ``run for \(N\) seconds''}

\textbf{Problem.}
In a windowed simulator, ``run for \(N\) simulated seconds'' cannot mean
arbitrary preemption at exactly that timestamp without violating scheduler
invariants.

\textbf{Resolution.}
Both interactive and socket controllers interpret bounded execution as
``continue until the first window boundary whose logical time is at least the
target''. This gives a precise, deterministic, and implementation-compatible
semantics to run-until control.

\section{Capabilities and Significance}

From a systems perspective, the contribution of this refactoring is not merely
that Shadow now accepts more commands. The deeper contribution is that Shadow
gains a structured \emph{temporal control plane} with clear semantics and
multiple front-ends, and that this control plane is connected to a
\emph{stateful restore path}. Together, these properties support workflows
that are central to modern systems debugging and testing:

\begin{itemize}[leftmargin=1.7em]
  \item \textbf{Time-travel debugging}: checkpoint before a fault, attach to
  the relevant native processes, and repeatedly restore for inspection.
  \item \textbf{Snapshot-based fuzzing}: branch from a promising temporal cut
  instead of replaying a long prefix for every trial.
  \item \textbf{Deterministic replay around fault windows}: reproduce
  distributed behavior around specific logical-time intervals with precise
  temporal control.
  \item \textbf{Cross-layer orchestration}: coordinate the simulator with
  external persistent services using a unified control protocol.
\end{itemize}

In short, the current implementation turns Shadow from a simulator that
\emph{internally keeps track of time} into a system that \emph{externally
exposes time as a controllable resource}.

\section{Discussion}

The present implementation should be understood as a principled, boundary-aware
time-control substrate rather than merely a convenience wrapper. Its central
methodological decisions---scheduler-aligned intervention, explicit control
results, reified checkpoint state, and orchestration-aware restore---provide
the right abstraction boundary for future work on stronger replay, richer
branching, and higher-level tooling.

At the same time, the design is intentionally conservative where full
reconstruction is not yet guaranteed. Some transient task categories remain
hard to reconstruct generically, and restore semantics for complex mixed
internal/external states still depend on orchestration discipline. These are
not accidental limitations; they follow directly from the fact that Shadow is
a hybrid system combining simulator state, native process state, IPC state,
and external persistent state. The contribution of the current design is that
these boundaries are now explicit and therefore manageable.

\section{Conclusion}

This refactoring establishes time control as a core systems abstraction in
Shadow. The resulting methodology is centered on one idea: every temporal
operation---pause, step, run-until, checkpoint, and restore---must be defined
with respect to a boundary-aligned, globally consistent cut of the parallel
event-driven execution. Building such a mechanism is inherently non-trivial in
Shadow because time is distributed across scheduler state, host-local events,
native managed processes, shared-memory IPC, and external persistent state.

By decomposing the problem into a pluggable control plane, explicit
decision/result types, serializable checkpoint representations, CRIU-assisted
process handling, and orchestrator-aware restore semantics, the current system
provides a practical foundation for advanced temporal workflows. This is the
key outcome of the work: Shadow now offers not just faster simulated time, but
a programmable temporal substrate for debugging, replay, and systematic
exploration.

\end{document}
